\hypertarget{user-guide---rnra-data-processing-and-error-management}{%
\section{USER GUIDE - RNRA data processing and error
management}\label{user-guide---rnra-data-processing-and-error-management}}

\hypertarget{introduction}{%
\subsection{1. Introduction}\label{introduction}}

\begin{itemize}
\tightlist
\item
  objectif du logiciel.
\item
  publique cible
\end{itemize}

\hypertarget{scientific-principle}{%
\subsection{2. scientific principle}\label{scientific-principle}}

Explication : - RNRA - Methode ANOVA - Hypothese statistique - Test de
normalité et homogénité

\hypertarget{detailed-installation}{%
\subsection{3. detailed installation}\label{detailed-installation}}

Procedure etape par etape

\hypertarget{interface-description}{%
\subsection{4. Interface description}\label{interface-description}}

\hypertarget{conversion}{%
\subsubsection{4.1 Conversion}\label{conversion}}

\hypertarget{interface-description-1}{%
\subsection{4. Interface Description}\label{interface-description-1}}

\hypertarget{conversion-window}{%
\subsubsection{4.1 Conversion Window}\label{conversion-window}}

This window allows the batch conversion of multiple \texttt{.mpa} files
contained in a selected folder.

The software processes all \texttt{.mpa} files in the chosen directory
and converts them into structured \texttt{.txt} files.

The conversion includes: - Extraction of raw acquisition data - Dead
time calculation - Separation of data for each ADC channel -
Initialization and formatting of the output data

\hypertarget{user-actions}{%
\paragraph{User Actions}\label{user-actions}}

\begin{enumerate}
\def\labelenumi{\arabic{enumi}.}
\tightlist
\item
  Select the input folder

  \begin{itemize}
  \tightlist
  \item
    The folder must contain the \texttt{.mpa} files generated by the MPA
    acquisition software.\\
  \item
    The path can be entered manually or selected using the
    \textbf{``Browse''} button.
  \end{itemize}
\item
  Select the output folder (optional)

  \begin{itemize}
  \tightlist
  \item
    The user may choose where the converted \texttt{.txt} files will be
    saved.\\
  \item
    If no output path is specified, the converted files will be stored
    in a temporary working directory for the duration of the current
    session. {TODO: Check this part later}
  \item
    Each \texttt{.mpa} file will generate one \texttt{.txt} file per ADC
    channel following the naming convention described above.
  \end{itemize}
\end{enumerate}

Once the input folder is defined, the batch conversion process can be
launched.

\hypertarget{output-text-file-structur}{%
\paragraph{Output Text File Structur}\label{output-text-file-structur}}

For each \texttt{.mpa} file contained in the selected folder, the
software generates one output \texttt{.txt} file per ADC channel.

The naming convention is:

filename\_adcname.txt

\textbf{File naming:} \texttt{filename\_ADCname.txt} - \texttt{filename}
: name of the original \texttt{.mpa} file - \texttt{ADCname} : ADC
channel identifier

Example:

If the input file is: sample01.mpa

And it contains ADC1 and ADC2,

The generated files will be: sample01\_ADC1.txt sample01\_ADC2.txt

\textbf{File content / Structure:} 1. \textbf{Header line:} - Indicates
the calculated dead time factor for this ADC channel.

\begin{enumerate}
\def\labelenumi{\arabic{enumi}.}
\setcounter{enumi}{1}
\tightlist
\item
  \textbf{Data lines:}\\
\end{enumerate}

\begin{itemize}
\tightlist
\item
  \texttt{Channel}: ADC channel number or identifier\\
\item
  \texttt{Count}: measured counts for each sample
\end{itemize}

\textbf{Notes:} - Each \texttt{.mpa} file generates one \texttt{.txt}
file per ADC channel.\\
- All files are formatted for direct analysis in Python, or other data
processing software.

\hypertarget{calibration}{%
\subsubsection{4.2 Calibration}\label{calibration}}

This window allows the user to perform a calibration from a
\texttt{.txt} file generated in the Conversion window.

The user can select a text file, then for each chosen peak(s):

\begin{enumerate}
\def\labelenumi{\arabic{enumi}.}
\tightlist
\item
  Identify the approximate position of the peak by clicking on it.
\item
  Fit a Gaussian function to the selected peak to determine its
  centroid.
\item
  Calculate a linear regression between the ADC channel numbers and the
  corresponding energy values of the chosen peaks.
\end{enumerate}

The result is a calibration curve that maps channels to energy values
for subsequent analyses.

\hypertarget{user-actions-1}{%
\paragraph{User Actions}\label{user-actions-1}}

\hypertarget{output-file-naming-convention}{%
\paragraph{Output File Naming
convention}\label{output-file-naming-convention}}

\hypertarget{processing}{%
\subsubsection{4.3 Processing}\label{processing}}

This window allows the user to perform all the steps of the data
processing workflow.

\begin{enumerate}
\def\labelenumi{\arabic{enumi}.}
\tightlist
\item
  \textbf{Loop over ROI (Region of Interest):}

  \begin{itemize}
  \tightlist
  \item
    For a given ROI in energy, the program counts the number of counts
    and normalizes it by the integrated charge.\\
  \item
    This produces the excitation profile for the selected ROI.
  \end{itemize}
\item
  \textbf{Eject peak (optional):}

  \begin{itemize}
  \tightlist
  \item
    For a given energy, the software searches for peaks in the output of
    the loop within a specified interval (ΔE = 200 keV by default).\\
  \item
    Points corresponding to peaks within another specified interval fix
    (20 keV by default) can be removed if necessary.\\
    {Can be change easily do an explenation how with change in the code
    and see if it works}
  \end{itemize}
\item
  \textbf{Sigmoid fit:}

  \begin{itemize}
  \tightlist
  \item
    Once the loop is complete (and peaks optionally ejected), the
    software performs a sigmoid fit on the resulting data.\\
  \item
    It calculates the difference between the high and low plateaus and
    assigns a group to each profile.
  \end{itemize}
\end{enumerate}

\hypertarget{user-actions-2}{%
\paragraph{User Actions}\label{user-actions-2}}

\hypertarget{output-file-naming-convention-1}{%
\paragraph{Output File Naming
convention}\label{output-file-naming-convention-1}}

\hypertarget{anova}{%
\subsubsection{4.4 ANOVA}\label{anova}}
